\documentclass[10pt,a4paper]{amsart}
\usepackage[margin=19mm]{geometry}
\usepackage{amsmath,amssymb,mathtools}
\usepackage[scr=rsfs,cal=euler]{mathalfa}
\usepackage{xfrac}
\usepackage[unicode,hidelinks,bookmarks=false]{hyperref}
\newcommand{\ie}{i.e.\ }
\newcommand{\cf}{cf.\ }
\newcommand{\resp}{respectively\ }
\newcommand{\Subsection}{\subsection}
\providecommand{\nobreakdash}{\nobreak\hbox{-}}
\setlength{\emergencystretch}{2em}
\multlinegap0pt
\usepackage{siunitx}
\usepackage{siunitx}
\usepackage{siunitx}
\usepackage{siunitx}
\usepackage{amssymb}
\usepackage{algorithm}\usepackage{algpseudocode}
\usepackage{bm}
\usepackage{xcolor}
\usepackage{siunitx}
\usepackage{float}
\newtheorem{assumption}{Assumption}
\newtheorem{lemma}{Lemma}
\newcommand{\R}{\mathbb{R}}
\newcommand{\be}{\begin{equation}}
\newcommand{\ee}{\end{equation}}
\title{Mirror Descent Linearized Augmented Lagrangian Methods for Nonconvex Constrained Stochastic Zeroth-Order Optimization}
\author{Qiankun Shi, Han Yuan, Xiao Wang, Hao Wang}
\date{Source: arXiv:2504.09409v3 (CC BY 4.0)}
\begin{document}
\maketitle

\noindent\emph{Source and licence.} Mirror Descent Linearized Augmented Lagrangian Methods for Nonconvex Constrained Stochastic Zeroth-Order Optimization. Version arXiv:2504.09409v3 (CC BY 4.0), distributed by arXiv under the Creative Commons Attribution 4.0 International licence, http://creativecommons.org/licenses/by/4.0/, as recorded in the arXiv licence metadata for this exact version and shown on the abstract page https://arxiv.org/abs/2504.09409v3. Full licence notice: \texttt{licenses/arxiv-2504.09409-CC-BY-4.0.txt}.\par\smallskip

\noindent\textit{Editorial scope.} Counted unit: the article's lemma lem:smooth-mirror-map (source0.tex lines 1929--1934). The article's complete original proof of it is delivered with it (source0.tex lines 1935--2059, one \texttt{proof} environment of 125 lines). No mathematics is written, repaired or re-derived: the statement and the proof are the article's own lines, in the article's own notation, and no retained line is substituted or re-flowed.  Retained uncounted as the source-local context the proof needs: lines 256--265; lines 433--443; lines 1188--1190; lines 1193--1197.  Nothing was omitted from any retained span, so the package carries no omission record.  No retained line contains a citation, so the wrapper carries no bibliography and no bibliography entry.  No retained line contains a figure environment, an image-inclusion command or drawing code, so no image is delivered and the package declares no asset dependency.  Editorial formatting only, in addition: an AMS \texttt{amsart} preamble that restates the article's own declarations and macros used by the retained lines, namely the environments \texttt{assumption}, \texttt{lemma} on separate counters, together with the standard packages those lines need.  The retained environments print in this document as Assumption 1, Lemma 1 in order of appearance, whereas the article prints the same items with the numbering of its own full document.  The complete native source of the article as distributed in the arXiv e-print for this exact version is bundled unchanged as \texttt{sources/176436/source0.tex}.
\begin{assumption}\label{ass:lip}
    The set $X \subset \R^d$ is nonempty, closed, convex and bounded w.r.t. the $q$-norm, i.e., there exists  $D_X>0$ such that $\|x\|_q \le D_X$, for any $x\in X$. The function $f:\R^d \rightarrow \R$ is continuously differentiable over $X$, and  $$\|\nabla f(x)-\nabla f(y)\|_2\le L_f\|x-y\|_2,\quad | f(x) - f(y) | \leq M_f\| x - y\|_2, \quad \forall  x,y\in \mathcal O_\nu,$$
    where $L_f, M_f>0$. 
    Function {$h:\R^d \to \R\cup \{\infty\}$} is proper, lower semicontinuous and convex over $X$, and 
    \[\sup_{g_h \in \partial h(x)} \|g_h\|_2 \leq M_h, \quad  \forall x\in X,\]
    where $M_h>0.$ 
    Functions $c_i:\R^d \rightarrow \R$, $i=1,\ldots,m$ are continuously differentiable over $X$, and 
    $$ |c_i(x)|\le \bar M_c, \quad \|\nabla c_i(x)\|_2\le M_c, \quad \|\nabla c_i(x)-\nabla c_i(y)\|\le L_c\|x-y\|, \quad \forall i=1,\ldots,m;\ x,y\in X,$$
    where $\bar M_c, M_c, L_c>0$. The objective function $\Phi(x)$ has a finite lower bound $\Phi_{\rm inf}$ over $X.$
\end{assumption}

\begin{assumption}\label{ass:v}
    The distance-generating function $v: \mathbb{R}^d \to \mathbb{R}$ is $1$-strongly convex   w.r.t. the \(\ell_q\)-norm, meaning 
\be\label{v_str_cvx}
\langle x - y, \nabla v(x) - \nabla v(y) \rangle \geq \|x - y\|_q^2, \quad  x,y\in \mathbb R^d.
\ee
And $v$ is $L_v$-smooth over $X$ w.r.t. the $\ell_q$-norm, that is
\be\label{v_Lv_sm}
\|\nabla v(x)-\nabla v(y)\|_p
    \le L_v\|x-y\|_q,\quad  x,y \in X. 
\ee
\end{assumption}

\begin{equation}\label{eq:smoothed_v}
v_\delta(x) = \frac{1}{2(q-1)} ( \sum_{i=1}^d (x_i^2 + \delta^2)^{q/2} )^{2/q},
\end{equation}

\begin{equation}\label{eq:Lv}
    L_{v}= \frac{2-q}{q-1} +
    \frac{1}{q-1}
    (D_X^q+d\delta^q)^{(2-q)/q}\delta^{q-2}.
\end{equation}

\begin{lemma}\label{lem:smooth-mirror-map}
Let $v_\delta$ be defined by \eqref{eq:smoothed_v}. %, where $\frac1p+\frac1q=1$.
Then $v_\delta$ satisfies Assumption \ref{ass:v}
%is 1-strongly convex w.r.t. the $\ell_q$-norm, and $\nabla v_\delta$ is $L_{v}$-Lipschitz continuous on $X$, where
with $L_v$ defined in \eqref{eq:Lv}.
\end{lemma}

\begin{proof}
By Assumption \ref{ass:lip}, $X$ is bounded in the $\ell_q$-norm with $D_X$.
Let
\( \Psi_\delta(x)=\sum_{i=1}^d (x_i^2+\delta^2)^{q/2},\) and \(a_i=x_i^2+\delta^2 .\)
Then, we obtain
\(
    v_\delta(x)=\frac{1}{2(q-1)}\Psi_\delta(x)^{2/q}.
\) 
A direct calculation gives the gradient entries
\[
    \nabla_i v_\delta(x)=\frac{1}{q-1}\Psi_\delta(x)^{2/q-1}a_i^{q/2-1}x_i ,\quad i=1,\ldots,d.
\]
Moreover, the Hessian entries are
\[
    [\nabla^2 v_\delta(x)]_{ij}
    =\frac{2-q}{q-1}\Psi_\delta(x)^{2/q-2}a_i^{q/2-1}x_ia_j^{q/2-1}x_j                +\frac{1}{q-1}\Psi_\delta(x)^{2/q-1}a_i^{q/2-2}\big((q-1)x_i^2+\delta^2\big)\mathbf 1_{{i=j}} 
\]
for $i,j=1,\ldots,d.$ 
Thus, we obtain\(\nabla^2 v_\delta(x)=M(x)+\Lambda(x),\) where \(M(x)=\frac{2-q}{q-1}\Psi_\delta(x)^{2/q-2}ww^\top, w_i=a_i^{q/2-1}x_i,\) and \(\Lambda_{ii}(x)=\frac{1}{q-1}\Psi_\delta(x)^{2/q-1}a_i^{q/2-2}((q-1)x_i^2+\delta^2), \Lambda_{ij}(x)=0, \) for \(i,j=1,\ldots,d;\ i\ne j .\)

We first prove 1-strong convexity. Since $M(x)$ is positive semidefinite, for any $y\in\mathbb R^d$,
\[
    y^\top \nabla^2 v_\delta(x)y
    \ge
    y^\top \Lambda(x)y     =
    \frac{1}{q-1}
\Psi_\delta(x)^{(2-q)/q}
    \sum_{i=1}^d
    a_i^{q/2-2}
    \big((q-1)x_i^2+\delta^2\big)y_i^2    
\ge\Psi_\delta(x)^{(2-q)/q}
    \sum_{i=1}^d
    a_i^{(q-2)/2}y_i^2 .
\]
The last inequality follows from\(\frac{(q-1)x_i^2+\delta^2}{q-1}=x_i^2+\frac{\delta^2}{q-1}\ge x_i^2+\delta^2=a_i .\) On the other hand, by H\"older's inequality, it holds
\[
    \sum_{i=1}^d |y_i|^q
    =
    \sum_{i=1}^d
    (a_i^{(q-2)/2}y_i^2)^{q/2}
    (a_i^{q/2})^{(2-q)/2}       
    \le
    (\sum_{i=1}^d a_i^{(q-2)/2}y_i^2)^{q/2}
    (\sum_{i=1}^d a_i^{q/2})^{(2-q)/2}.
\]
Increase both sides to the power $2/q$ yields \(\|y\|_q^2\le\Psi_\delta(x)^{(2-q)/q}\sum_{i=1}^d a_i^{(q-2)/2}y_i^2 .\) Consequently, we have \(y^\top \nabla^2 v_\delta(x)y\ge\|y\|_q^2 ,\ \forall x,y\in\mathbb R^d\), applying the fundamental theorem of calculus gives
\(\langle x-y,\nabla v_\delta(x)-\nabla v_\delta(y)\rangle=\int_0^1(x-y)^\top\nabla^2v_\delta(y+t(x-y))(x-y)\,dt\geq
\|x-y\|_q^2\), which proves that $v_\delta$ is 1-strongly convex w.r.t. the $\ell_q$-norm.

We next proof $v_\delta$ is \(L_v\)-smooth on $X$.
 H\"older's inequality gives
\(
    \langle w,h\rangle^2
    \le
    \|w\|_p^2\|h\|_q^2.
\) 
Since \(|w_i|=a_i^{q/2-1}|x_i|\le a_i^{(q-1)/2}\), we have \(\|w\|_p^p\le\sum_{i=1}^d a_i^{p(q-1)/2}=\sum_{i=1}^d a_i^{q/2}=\Psi_\delta(x),\) where $p=q/(q-1)$. Hence \(\|w\|_p^2\le\Psi_\delta(x)^{2/p}=\Psi_\delta(x)^{2(q-1)/q}.\) Therefore,
\[
    y^\top M(x)y
    =
    \frac{2-q}{q-1}
 \Psi_\delta(x)^{2/q-2}
    \langle w,y\rangle^2  
    \le
    \frac{2-q}{q-1}   \Psi_\delta(x)^{2/q-2}    \Psi_\delta(x)^{2(q-1)/q}
    \|y\|_q^2        =
    \frac{2-q}{q-1}\|y\|_q^2 .
\]
On the other hand, since $q\in(1,2]$, we have \(\frac{(q-1)x_i^2+\delta^2}{x_i^2+\delta^2}\le 1,\) and \(a_i^{(q-2)/2}\le\delta^{q-2},\) thus
\[
    \Lambda_{ii}(x)
    =
    \frac{1}{q-1}
    \Psi_\delta(x)^{(2-q)/q}
    a_i^{(q-2)/2}
    \frac{(q-1)x_i^2+\delta^2}{x_i^2+\delta^2}              \le
    \frac{1}{q-1}
    \Psi_\delta(x)^{(2-q)/q}
    \delta^{q-2}.
\]
Moreover, for $x\in X$, we have \(\Psi_\delta(x)=\sum_{i=1}^d (x_i^2+\delta^2)^{q/2}\le\sum_{i=1}^d |x_i|^q+d\delta^q\le D_X^q+d\delta^q ,\) which derives
\(    \Lambda_{ii}(x)
    \le
    \frac{1}{q-1}
    (D_X^q+d\delta^q)^{(2-q)/q}\delta^{q-2}, \forall i=1,\ldots,d\) and \(x\in X.
\) 
Since $\|y\|_2\le \|y\|_q$ for $q\in(1,2]$, we obtain \[y^\top \Lambda(x)y\le\frac{1}{q-1}(D_X^q+d\delta^q)^{(2-q)/q}\delta^{q-2}\|y\|_q^2 .\]
Combining the bounds for $y^\top M(x)y$ and $y^\top \Lambda(x)y$, we have
\(
    y^\top \nabla^2 v_\delta(x)y
    \le
    L_{v}\|y\|_q^2,
    \) \(\forall x\in X,  y\in\mathbb R^d,
\)
where $L_v$ is given by \eqref{eq:Lv}.
We next show that this quadratic-form bound implies the corresponding
operator-norm bound. Since $\nabla^2 v_\delta(x)$ is symmetric positive
semidefinite, the Cauchy--Schwarz inequality w.r.t. the
semidefinite bilinear form induced by $\nabla^2 v_\delta(x)$ gives, for
any $y,z\in\mathbb{R}^d$,
\[
    |z^\top \nabla^2 v_\delta(x) y|
    \le
    \sqrt{z^\top \nabla^2 v_\delta(x) z}\,
    \sqrt{y^\top \nabla^2 v_\delta(x) y} 
    \le
    L_v \|z\|_q \|y\|_q .
\]
Then by the duality between the $\ell_q$ and $\ell_p$ norms,
\(\|\nabla^2 v_\delta(x)y\|_p
    = \sup_{\|z\|_q\le 1}
       |z^\top \nabla^2 v_\delta(x)y| 
    \le L_v\|y\|_q .\)
Hence,
\[
    \|\nabla^2 v_\delta(x)\|_p
    =\sup_{\|y\|_q\le1}
      \|\nabla^2 v_\delta(x)y\|_p
    \le L_v,
    \qquad \forall x\in X.
\]
Since $X$ is convex, for any $x,y\in X$, the segment $x+t(y-x)$ remains in $X$. Hence,
\[\|\nabla v_\delta(y)-\nabla v_\delta(x)\|_p\le\int_0^1\|\nabla^2v_\delta(x+t(y-x))(y-x)\|_p\,dt \le L_{v}\|y-x\|_q .\]
Therefore, $v_\delta$ is $L_{v}$-smooth on $X$ w.r.t. the $\|\cdot\|_q$-norm, which completes the proof.
\end{proof}

\end{document}
